% Counting and parity: Z11 (C8), Z16 (C7).
% Open questions for the author:
% - Z11's credit adds Arasu-Seberry 1998, Thm 1(iii)-(iv) (read by the refuter in R5, from OCR);
%   R5 notes that Wallis-Whiteman attribute the bound to their reference [2], not opened.
% - Z16: R12 notes that the coset facts |a_h| <= m_h <= c, m_h = a_h (mod 2) are also in He-Wu,
%   Lemma 4.1; the remark says so. solutions.md's Origin says only "ours".
% - MarkedCounts.lean's doc-string says weights exist "whenever no real counts ... meet both
%   bounds". By Farkas' lemma a one-sided certificate exists only when the relaxation with that
%   one bound is infeasible; the text below says this. Lean is unaffected (it checks weights).
% - The examples (N318, N1767) were checked by hand and by a script against the Lean certificates.

\section{Counting and parity}\label{sec:counting}

The two criteria of this section count \emph{marked} elements, those $x$ with $D(x) \ne 0$:
on $G$ itself (Theorem~\ref{thm:z11}), and coset by coset on a quotient
(Theorem~\ref{thm:z16}).

\subsection{Support parity}

\begin{theorem}\label{thm:z11}
Let $G$ be a finite abelian group of odd order $v$, let $k \ge 0$ and $\lambda$ be integers with
$\lambda$ even, and put $d = \max(v - k, 0)$. If $d(d - 1) < v - 1$, then $G$ has no SDS with
parameters $(v,k,\lambda)$.
\end{theorem}

\begin{proof}
Let $D$ be such an SDS and $Z = \{x : D(x) = 0\}$ its set of unmarked elements; $|Z| = v - k = d$.
Fix $g \ne 0$. Modulo~$2$, $D(x) D(x+g)$ is $1$ exactly when $x$ and $x + g$ are both marked. By
inclusion--exclusion the number of such $x$ is $v - 2d + |Z \cap (Z - g)|$, so
\[
  \lambda \equiv v - 2d + |Z \cap (Z - g)| \pmod 2 .
\]
With $v$ odd and $\lambda$ even, $|Z \cap (Z - g)|$ is odd, in particular nonzero: $g = z' - z$
for some $z, z' \in Z$. So the $d(d-1)$ ordered pairs of distinct elements of $Z$ give all
$v - 1$ nonzero differences, and $d(d - 1) \ge v - 1$.
\end{proof}

The bound is sharp: the circulant weighing matrix $(7, 4, 0)$ and a $(13, 9, 6)$ SDS have
$d(d-1) = v - 1$. In that case every $g \ne 0$ is a difference of unmarked elements in exactly
one way, so $Z$ is a planar difference set.

\begin{example}
The case N318, $(47, 41, 20)$ in $C_{47}$, has $v$ odd, $\lambda$ even and $d = 6$, and
$d(d-1) = 30 < 46$.
\end{example}

\begin{remark}
For $\lambda = 0$ this is the classical bound for weighing matrices of odd
order~\cite{wallis-whiteman}, \cite[Thm.~1(iii)--(iv)]{arasu-seberry}, the latter for circulant
ones, with the planar difference set in the case of equality. For general $\lambda$ it is ours;
He and Wu~\cite[Cor.~4.5]{he-wu} cover $d = 1$ with $v = p^m$, $p \equiv 3 \pmod 4$, $m$ even.
The criterion settles \zrows{11} cases. Lean:
\lean{SignedDifferenceSets.not_hasSDS_of_supportParity}; the cases are checked in the module
\lean{SignedDifferenceSets.ParityRows}.
\end{remark}

\subsection{Marked elements per coset}
Let $\phi \colon G \to H$ be surjective with fibres of size $c$, and let $m_h$ be the number of
marked elements in the fibre over $h$. For an integer $x$ let $M_c(x)$ be $c$ if
$x \equiv c \pmod 2$ and $c - 1$ otherwise: the largest number of marked elements a fibre with
sum $x$ can hold.

\begin{lemma}\label{lem:marked}
For an SDS with parameters $(v,k,\lambda)$ and coset sums $a$ on $H$,
\[
  |a_h| \le m_h \le M_c(a_h) \quad (h \in H), \qquad \sum_h m_h = k, \qquad
  \text{hence} \quad \sum_h |a_h| \le k \le \sum_h M_c(a_h).
\]
\end{lemma}

\begin{proof}
$a_h$ is a sum of $m_h$ signs, so $|a_h| \le m_h$ and $m_h \equiv a_h \pmod 2$; and
$m_h \le c$, which with the parity gives $m_h \le M_c(a_h)$. The fibres partition $G$, and $D$
has $k$ marked elements.
\end{proof}

We keep the notation of Section~\ref{sec:congruences}: $N = |H|$, $Q = k + \lambda(c - 1)$,
modulus $q$, residue $t_0$, and the box test (Definition~\ref{def:box}). Let
$V(t_0) = \{x \in \Z : |x| \le c,\ x \equiv t_0 \pmod q\}$. For integers $\alpha, \beta, \gamma$
put $P(x) = \gamma + \beta x + \alpha x^2$ and $T(s') = \gamma N + \beta s' + \alpha Q$, which is
$\sum_h P(a_h)$ when $\sum_h a_h = s'$ and $\sum_h a_h^2 = Q$.

\begin{theorem}\label{thm:z16}
Let $G$, $k \neq \lambda$, $\phi \colon G \to H$, $c$, $p, i, j$ and $s$ be as in
Theorem~\ref{thm:z9}, and $q = p^j$. Suppose that for both $s' \in \{s, -s\}$ and every
$t_0 \in \{0, \dots, q - 1\}$, either $(N, q, s', Q, t_0)$ fails the box test with $c$, or there
are integers $\alpha, \beta, \gamma$ and $d \ge 0$ such that
\begin{enumerate}
\item[(L)] $P(x) \le d\,|x|$ for all $x \in V(t_0)$, and $T(s') > dk$; or
\item[(U)] $P(x) \ge d\,M_c(x)$ for all $x \in V(t_0)$, and $T(s') < dk$.
\end{enumerate}
Then $G$ has no SDS with parameters $(v,k,\lambda)$.
\end{theorem}

\begin{proof}
Let $D$ be such an SDS, $s' = \sum_h a_h \in \{s, -s\}$, and $t_0$ the residue of $a_0$ modulo
$q$. By Lemma~\ref{lem:sc-congr} and Lemma~\ref{lem:contraction}(1), every $a_h$ lies in
$V(t_0)$; they are not all equal (proof of Theorem~\ref{thm:z8}), so by Lemma~\ref{lem:box}
the box test passes. So there are weights as in (L)
or~(U), and $\sum_h P(a_h) = T(s')$. In case (L), Lemma~\ref{lem:marked} gives
$T(s') = \sum_h P(a_h) \le d \sum_h |a_h| \le dk$. In case (U) it gives
$dk = d \sum_h m_h \le d \sum_h M_c(a_h) \le \sum_h P(a_h) = T(s')$. Both contradict the
hypothesis.
\end{proof}

The weights are a dual (Farkas) certificate. Let $n_x$ be the number of $h$ with $a_h = x$, for
$x \in V(t_0)$. Then $n_x \ge 0$, $\sum_x n_x = N$, $\sum_x x\, n_x = s'$,
$\sum_x x^2 n_x = Q$, and, by Lemma~\ref{lem:marked}, $\sum_x |x|\, n_x \le k$ and
$\sum_x M_c(x)\, n_x \ge k$. By Farkas' lemma, weights as in (L) exist exactly when this system
without its last inequality has no real solution, and weights as in (U) exactly when it has
none without the inequality before; the weights can be taken integral. We find them by linear
programming; Lean only checks the finitely many inequalities on $V(t_0)$, and needs only the
direction proved above.

\begin{example}
The case N1767, $(95, 84, 40)$ in $C_{95}$, has $n = 44$ and $s = 62$. Take $H = C_{19}$
($c = 5$), $p = 2$, $j = 1$, $i = 9$ ($2^9 \equiv -1 \pmod{19}$). For $t_0 = 1$ the nineteen
coset sums would be odd, with even sum $\pm 62$; the box test fails. For $t_0 = 0$ it passes
for both signs. But then every coset sum is even, so every fibre of size $5$ holds an even
number, at most $4$, of marked elements, and $k \le 19 \cdot 4 = 76 < 84$. This is (U) with
$(\alpha, \beta, \gamma, d) = (0, 0, 4, 1)$: $M_5(x) = 4$ on $V(0) = \{0, \pm 2, \pm 4\}$, and
$T(\pm 62) = 76 < 84$.
\end{example}

\begin{remark}
The facts of Lemma~\ref{lem:marked} on a single coset are also in~\cite[Lemma~4.1]{he-wu}; the
criterion and its certificates are ours. It settles \zrows{16} cases. Lean:
\lean{SignedDifferenceSets.not_hasSDS_of_marked}; the cases are checked in the modules
\lean{SignedDifferenceSets.MarkedCountsRows.Part01} to \lean{Part03}.
\end{remark}
